% =====================================================================
% I. INTRODUCTION -- OWNER: Patrick (first draft 2026-10-20, revised 2026-10-22, written
% with Claude Code)
% Sources: journal/2026-09-25.md; readings/jacovi_goldberg_2020.md;
% readings/lanham_2023.md; readings/parcalabescu_2024.md; journal/2026-10-15.md;
% journal/2026-10-20.md; journal/2026-10-21.md; outputs/clean/paper_stats.json;
% outputs/clean/copy_vs_use_summary.json; outputs/raw/copy_vs_use.config.json (0.6B).
% Claims about turpin2023, chen2025, wei2022 and qwen3 follow their abstracts.
% =====================================================================
\section{Introduction}\label{sec:intro}
Reasoning models write out their working before answering~\cite{wei2022}, and Qwen3
offers thinking and non-thinking modes in one model~\cite{qwen3}. Readers may take the
working as an explanation, but an explanation can be plausible without being faithful
to the process behind the prediction~\cite{jacovi2020}. Most faithfulness tests
intervene, corrupting the reasoning~\cite{lanham2023} or adding hints to the
prompt~\cite{turpin2023,chen2025}. We first ask a narrower, label-free question of
unedited output: is the final answer \emph{consistent} with the written arithmetic?
Because consistency is weaker than faithfulness~\cite{parcalabescu2024}, we pair it
with two interventions. Our contributions are:
\begin{itemize}
\item an exact-arithmetic consistency checker, and a Z3 logic
  checker~\cite{demoura2008} validated on FOLIO annotations~\cite{han2024folio};
\item zero answer--working contradictions in 83 comparable Qwen3-1.7B responses on
  GSM8K~\cite{cobbe2021};
\item a copy-vs-use control on Qwen3-0.6B, indicating that first-digit sensitivity to
  planted working mostly reflects copying rather than computation;
\item a misleading-hint test in which every hint-following answer acknowledged the
  hint, but the hint left previously correct responses unfinished.
\end{itemize}
